% =====================================================================
%  IV. EXPERIMENTATION
%  Sources: notes/experimental-protocol.md S4, S5, S8, S9, S11; gate 1
%  result (protocol S6.2, PR #4); gate 2 result (protocol S10.1, PR #5).
%  Numbers in the two gate subsections are measured. Everything about
%  labelling and routing describes the plan and must be re-checked against
%  what was actually run.
% =====================================================================
\section{Experimentation}
\label{sec:experimentation}

\subsection{Hardware and Software}
\label{sec:hardware}

All local experiments run on one laptop: an AMD Ryzen~7 8845HS CPU with
16\,GB of system memory and an NVIDIA GeForce RTX~4060 Laptop GPU with
8\,GB of memory, driver 616.56 and an enforced board power limit of 115\,W.
Models are served by Ollama 0.34.0~\cite{ollama} from GGUF weights
quantised to Q4\_K\_M~\cite{gerganov2023llamacpp}.
Every model is addressed by an explicit quantisation tag, and the weight
digest, parameter count, context length and chat template of each are
recorded in a manifest before any generation, so that a silently updated
tag cannot change the models between runs. All four models report native
thinking support and a 40{,}960-token context. The 8\,GB memory limit means
only one model is loaded at a time. Runs are therefore organised
\emph{model-major}: each model is loaded once, all prompts are run in both
modes, and the model is unloaded before the next is loaded. Within a model,
prompts are processed in batches of 20, and each batch runs its two modes as
separate measurement windows. Which mode goes first alternates from batch to
batch, so that slow drift in temperature or background load cannot line up
with one mode. Other GPU workloads share the machine, so the harness waits
for 30\,s of continuous GPU idleness before loading a model or repeating a
window, and it checks that the whole model is placed on the GPU. The
software environment is Python 3.11, pinned with a lock file.

\subsection{Decoding}
\label{sec:decoding}

Greedy decoding is discouraged for Qwen3's thinking mode, so each mode uses
its recommended sampling settings~\cite{qwen2025modelcard}: temperature 0.6, top-$p$ 0.95 and top-$k$
20 for thinking, and temperature 0.7, top-$p$ 0.8 and top-$k$ 20 for
non-thinking, with min-$p$ 0 in both. Every generation receives its own seed,
derived deterministically from the prompt, cell, sample index and run
variant, and the seed is stored with the output.

The labelling runs draw one sample per prompt and cell ($k = 1$). One sample
makes $s(q, c)$ a noisy estimate of the cell's chance of success on that
prompt, so a noise subset of 25 prompts per family receives two further
samples in every cell, and the disagreement between samples measures how
often a label would change on resampling. Three samples throughout were the
original plan, but gate 2 put one sample across the eight cells at about
112\,s per prompt (Section~\ref{sec:gate-mode}), and tripling that would have
tripled a labelling run that already takes days on this machine.

Generation length is capped, and an answer that hits the cap counts as a
failure, because a truncated answer is of no use in deployment. Thinking
mode is capped at 4{,}096 new tokens, trace included, and non-thinking mode
at 2{,}048. The caps matter more than we expected: gate 2 showed that
truncation causes most thinking-mode failures at 4{,}096 tokens
(Section~\ref{sec:gate-mode}), and a longer cap would lengthen the run most on
exactly the prompts where thinking struggles. The oracle label therefore
depends on the caps, and we measure by how much. A sample of capped
generations, up to 40 per size and mode and spread evenly over families, is
rerun with an 8{,}192-token cap and a 10{,}240-token context, and we report
how often the longer cap turns a failure into a success and how often that
success would move $\cstar(q)$ to a cheaper cell. Before measurement, each
Ollama server session
runs one thinking and one non-thinking warm-up generation, which are
discarded; the gate showed a one-off prefill stall of about 6\,s at the
start of a session that would otherwise contaminate the first measurements.

\subsection{Corpus}
\label{sec:corpus}

The corpus covers four task families (Table~\ref{tab:corpus}). Three can be
graded programmatically. The fourth, open-ended requests, is the traffic
that motivates the problem and needs a judge.

\begin{table}[!t]
\centering
\caption{Corpus composition. Each gradable family has 400 prompts, of which
100 form a nested pilot; prompt counts per source are in parentheses. The
open-ended family has not yet been sampled (Section~\ref{sec:judge}).}
\label{tab:corpus}
\renewcommand{\arraystretch}{1.1}
\begin{tabularx}{\columnwidth}{@{}l >{\raggedright\arraybackslash}X >{\raggedright\arraybackslash}X@{}}
\toprule
Family & Sources & Grading \\
\midrule
Math & GSM8K~\cite{cobbe2021gsm8k} (200), MATH-500~\cite{hendrycks2021math,lightman2024verify} (200) & Normalised final answer \\
Code & HumanEval~\cite{chen2021humaneval} (154), MBPP~\cite{austin2021mbpp} (246) & Unit tests in a sandbox \\
Multiple choice & MMLU~\cite{hendrycks2021mmlu} (200), ARC-Challenge~\cite{clark2018arc} (200) & Option letter \\
Open-ended & Dolly-15k~\cite{conover2023dolly} (pending) & Judge (Sec.~\ref{sec:judge}) \\
\bottomrule
\end{tabularx}
\end{table}

Prompts are drawn with a fixed, recorded seed. Each source is walked in a
round-robin order over strata of an a-priori difficulty proxy where the
source has one---MATH level, MMLU subject and MBPP task band---so that any
prefix of the order is balanced across strata. The pilot takes the first
rows of that order and is therefore nested in the full corpus, and the first
25 pilot prompts of each family form the noise subset. The 60 prompts used in
gate 2 are excluded. HumanEval then has 154 problems, so the code family takes
all of them and the remaining 246 from MBPP. Near-duplicates are removed by
hashing normalised prompt text, across sources as well as within them.
Prompts longer than 2{,}048 tokens under the Qwen3 tokenizer are excluded,
so that context length does not vary with the size axis; in the sampled
corpus the median prompt has 83 tokens and the longest 862. Each source
dataset's row count and a content fingerprint are stored with the sample.

After labelling, the corpus is split 60/20/20 into training, validation and
test sets by prompt, stratified by family and by $\cstar(q)$. The label used
for stratification is computed in tokens, so that every cost unit is
evaluated on the same split. Router hyperparameters and $\tau$ are chosen on
the validation split, and the test split is used once, at the end.

\subsection{Oracle Labelling and Success Criteria}
\label{sec:labelling}

Every prompt is run through all eight cells. For each prompt the labelling
run records the full success vector $s(q, c_1), \dots, s(q, c_8)$ and the
measured token counts, latency and energy of every cell. The label used for
the classification ablation is $\cstar(q)$ from \eqref{eq:cstar}, but the
routers are trained on the vector, which carries more information and is
what allows $\tau$ to be swept at test time.

Success is decided per family. For math, the final answer is extracted
(from a boxed expression when present, otherwise the last number), normalised
by removing thousands separators, units and \LaTeX{} wrappers, and compared
with the reference. For code, the generated function is executed against
the dataset's unit tests in a subprocess with a timeout and no network
access, and it succeeds only if every test passes. For multiple choice, the
chosen option is parsed from the answer; an unparseable answer is a failure,
not a reason to resample. Open-ended answers are judged as described next.

\subsection{Judging Open-Ended Answers}
\label{sec:judge}

A judge drawn from the routed pool would grade its own family, so the judge
is Llama~3.1 8B Instruct~\cite{grattafiori2024llama3} at Q4\_K\_M, run
through Ollama only after labelling has finished. The open-ended family is
labelled after the three gradable families, and results in
Section~\ref{sec:results} cover the gradable families unless stated
otherwise.

The judge sees the request, the human-written Dolly response as a reference
and the visible answer, and scores the answer from 1 to 5 on a fixed rubric;
a score of 4 or more is a success. For open requests such as brainstorming,
the rubric treats the reference as one acceptable answer among many and asks
for content, not length, to be judged. Answers are scored one at a time
rather than in pairs, in a seeded random order, to limit the position and
verbosity biases documented for model judges~\cite{zheng2023judging}.
Decoding is greedy and constrained to a JSON reply, so scores are
reproducible. An answer cut off by its token cap, or with no visible text,
fails without being judged.

The judge is validated from two sides before its labels are used. About 200
judged answers, stratified by Dolly category and verdict, are re-scored with
the same rubric by a frontier model accessed through the OpenAI API
\decision{OpenAI model, pending approval of the API spend}, and the first
author scores about 100 of them blind to model, mode and judge score.
Agreement is reported as Cohen's $\kappa$~\cite{cohen1960kappa} on the
success decision and as quadratic-weighted $\kappa$~\cite{cohen1968weighted}
on the scores. The rule for using the labels was fixed before any judge
output existed and rests on the judge--human $\kappa$ for the success
decision: at 0.60 or above, the open-ended family is reported with the
others; from 0.40 to 0.60, it carries an explicit caveat; below 0.40, it
moves to a secondary analysis. The cut-offs are the conventional boundaries
of moderate and substantial agreement~\cite{landis1977measurement}.

\subsection{Baselines}
\label{sec:baselines}

Routers are compared with: the eight static policies that always use one
cell; uniform random assignment; the oracle, which always picks $\cstar(q)$
and is unreachable in practice; the four manual escalation policies of
Table~\ref{tab:policies} under both detection regimes; and the reactive
cascade of Section~\ref{sec:cascade}.

\subsection{Metrics}
\label{sec:metrics}

\begin{itemize}
  \item \emph{Cost--quality frontier.} Solved rate against net mean cost as
        $\tau$ varies, in tokens, seconds and joules. The area under a
        frontier is the best solved rate achievable within a cost budget,
        averaged over budgets from the cheapest to the most expensive static
        policy, so it lies between 0 and 1 in every unit.
  \item \emph{Escalation tax.} $\Tax$, $\Esc$ and $\Stk$ from
        \eqref{eq:decomp}, per policy and per router, in absolute terms and
        relative to the cost of $\cstar(q)$.
  \item \emph{Net saving.} $S_{\mathrm{net}}$ from \eqref{eq:netsaving}
        against each baseline, with the gross saving beside it.
  \item \emph{Cost per solved query} in each unit, including joules per
        solved query as defined in Section~\ref{sec:energy}.
  \item \emph{Regret decomposition.} The over-routing rate and its cost
        premium, and the under-routing rate and the cost of recovering from
        it, reported separately and never combined into one accuracy figure.
  \item \emph{Router overhead.} Median and 95th-percentile decision latency
        and joules per decision. For the cascade, the same figures are
        reported per verification check.
\end{itemize}

\subsection{Statistical Analysis}
\label{sec:stats}

Every headline figure carries a 95\% percentile bootstrap confidence
interval from 10{,}000 resamples over prompts~\cite{efron1993bootstrap}. All policies and routers evaluated on the same
prompts share one matrix of resampled prompt indices, so every comparison is
paired by prompt and its interval is taken over the resampled differences.
Two-sided $p$-values are obtained by inverting the percentile interval, as
twice the smaller share of resampled differences on either side of zero,
with a floor of $10^{-4}$. The Holm--Bonferroni procedure~\cite{holm1979simple} corrects the $p$-values within each group of
comparisons that share a cost unit, recovery rule and detection regime. For
routers, $\tau$ is fixed on the validation split before the test split is
resampled, so the intervals do not include the uncertainty of that choice.
Each trainable router is trained with three seeds; intervals are reported
for the first seed and the spread across seeds beside them. Label noise is
estimated on the noise subset, from how often the three samples of a cell
disagree and how often $\cstar(q)$ changes when it is computed from a
different sample.

\subsection{Validation Gate 1: Power Telemetry}
\label{sec:gate-telemetry}

NVML does not expose board power on every laptop GPU, and a measured-energy
result is only as good as the sensor behind it, so the telemetry was tested
before anything else. The test recorded 10\,s idle before any CUDA context
was created, 8\,s of half-precision $2048 \times 2048$ matrix multiplication
at about 90\% utilisation, and 8\,s of cool-down, polling both NVML sources
every 100\,ms. Table~\ref{tab:gate1} gives the result.

\begin{table}[!t]
\centering
\caption{Gate 1: the two NVML energy sources on the RTX~4060 Laptop GPU
(enforced limit 115\,W). Counter values are the mean power implied by the
counter's increase over each window.}
\label{tab:gate1}
\renewcommand{\arraystretch}{1.1}
\begin{tabular}{@{}l r r r l@{}}
\toprule
Source & Idle & Load & Cool-down & Verdict \\
\midrule
Sampled power & 2.4\,W & 75.1\,W & 10.5\,W & Usable \\
Energy counter & 88.3\,W & 127.6\,W & 256.9\,W & Rejected \\
\bottomrule
\end{tabular}
\end{table}

Sampled power behaves as expected. It is low at idle, follows the step load
within about a second, and never exceeds the enforced limit; its maximum
under load was 88.4\,W. The cumulative energy counter is the more
instructive result. It is reported as supported and it increases
monotonically, so a routine check accepts it, and our own first-pass check
did. Yet the power it implies is 88\,W with the GPU idle, 128\,W under load,
which is above the enforced limit, and 257\,W during cool-down, which is
higher than under load. A units error would scale the three windows by one
factor and keep their order; these values do not follow real power draw at
all. Only the comparison with sampled power exposed the fault, with a 41\%
disagreement under load and a factor of 36 at idle. We therefore integrate
sampled power only, and the batch-level method of Section~\ref{sec:energy}
follows from the 0.5\,s refresh period measured in the same test. For other
studies on laptop GPUs, the practical lesson is that a supported, monotonic
counter should still be cross-checked against sampled power and against the
enforced limit before its readings are trusted.

\subsection{Validation Gate 2: Separation of Thinking Modes}
\label{sec:gate-mode}

The grid assumes that a model's thinking and non-thinking modes are two
different operating points. Wang \emph{et al.} found only partial separation
in hybrid-thinking models~\cite{wang2025demystifying}, so we tested the
assumption on our own checkpoints before labelling. The test used 60 prompts,
ten from each of GSM8K, the 500-problem MATH test subset~\cite{lightman2024verify}, HumanEval, MBPP, MMLU and
ARC-Challenge, and ran each prompt once per model size in both modes: 480
generations in 1\,h\,54\,min. Both modes had the same 4{,}096-token cap. An
asymmetric cap would have built the length difference into the test and
truncated precisely the long non-thinking answers that indicate leakage.

Three criteria were fixed before the run, and a size counted as separated
only if it met all three: an area under the ROC curve of at least 0.90 for
thinking-mode outputs being longer than non-thinking outputs of the same
prompt; at most 5\% of non-thinking responses containing a thinking trace
or a literal \texttt{\textless think\textgreater} tag; and a density of reasoning markers
(``wait'', ``hmm'', ``let me'', ``alternatively'' and eight similar phrases)
in non-thinking answers no more than half that in thinking traces.

\begin{table}[!t]
\centering
\caption{Gate 2: mode separation by model size (60 prompts per size).
Ratio is the median of per-prompt think/no-think token ratios with a 95\%
bootstrap interval. Marker ratio is no-think marker density divided by
think-trace marker density.}
\label{tab:gate2}
\renewcommand{\arraystretch}{1.1}
\setlength{\tabcolsep}{3.5pt}
\begin{tabular}{@{}l c c c c c@{}}
\toprule
Size & AUC & Tokens (think / no) & Ratio (95\% CI) & Leak & Marker \\
\midrule
0.6B & 0.964 & 1444 / 195 & 13.9 (9.4--19.7) & 0 & 0.040 \\
1.7B & 0.954 & 2053 / 184 & 9.1 (7.6--11.9)  & 0 & 0.034 \\
4B   & 0.928 & 1429 / 149 & 7.0 (5.3--11.5)  & 0 & 0.042 \\
8B   & 0.953 & 1814 / 180 & 9.7 (7.7--17.0)  & 0 & 0.045 \\
\bottomrule
\end{tabular}
\end{table}

All four sizes passed (Table~\ref{tab:gate2}). No non-thinking response
contained a thinking trace or tag, reasoning-marker density in non-thinking
answers was 3--5\% of that in thinking traces, and the thinking output was
longer for every prompt at three of the four sizes and for 59 of 60 at 4B
(Wilcoxon signed-rank $p < 2 \times 10^{-11}$ at every size). On these
checkpoints, the leakage described by Wang \emph{et al.} does not appear by
these measures. The measures have limits: marker counting is a lexical proxy
rather than a semantic judgement, and 60 prompts per size cannot rule out
leakage on task types we did not sample.

The run also exposed two properties of the grid that shape the labelling
design. First, thinking dominates the compute budget. Across the eight
cells, one sample per prompt took 112\,s on average, and the four thinking
cells accounted for 88\% of that time. Second, the cap turns long thinking
into failure. Between 22\% and 32\% of thinking responses hit the
4{,}096-token cap, depending on size, and every capped response on the
graded families was wrong. On
the math and multiple-choice prompts, capped responses made up most
thinking-mode errors from 1.7B upward, and thinking-mode accuracy fell below
non-thinking accuracy at 1.7B, 4B and 8B; with capped responses excluded,
thinking was within a point of non-thinking at 8B and ahead at 1.7B. These
figures come from 40 graded responses per mode and size and approximate
grading, so they motivate a design check rather than support a claim: the
oracle label depends on the cap, and the labelling runs therefore include
the cap-sensitivity reruns described in Section~\ref{sec:decoding}. The
thinking cap stays at 4{,}096 tokens. The non-thinking cap is set at
2{,}048, more than ten times the median non-thinking answer length in
Table~\ref{tab:gate2}.
% Accuracy detail (gate 2, math+MCQ, think / think-uncapped / no-think):
% 0.6B 0.42/0.52/0.33; 1.7B 0.57/0.79/0.70; 4B 0.68/0.84/0.90;
% 8B 0.78/0.94/0.95. At 4B, uncapped think (0.84) is still below no-think.

\subsection{Public-Benchmark Arm}
\label{sec:public-arm}

The local pool stops at 8B parameters, and frontier tiers may behave
differently. To test whether the escalation-tax analysis carries over, the
same policy and tax computations are applied, without local inference, to
the per-model outcomes and costs published with
LLMRouterBench~\cite{li2026llmrouterbench}. Each of its records stores a
score, prompt and completion token counts, a dollar cost and a latency for
one model on one query, which is what the resubmission term needs: a
re-asked query pays its prompt tokens again at the stronger model's price.
RouterBench~\cite{hu2024routerbench} carries similar fields, but it is
distributed as serialised Python objects, which execute code when loaded,
so we do not use it. Energy is not available for these pools, so this arm
reports tokens, dollars and latency only. \tbd{public-benchmark arm: run
once the data download is approved}

\subsection{Reproducibility}
\label{sec:repro}

Each experiment is driven by a configuration file, and no experimental
constant lives in code. Seeds are set and stored for every generation and
every training run. Each run writes a manifest with the code commit,
configuration, package versions, driver version, GPU model and model
digests. Raw generations are stored with their token counts, timings, seeds
and energy. The processed tables that every result is computed from contain
prompt identifiers but no prompt text, and the corpus can be rebuilt from
the public source datasets with the recorded seed and fingerprints.
\decision{code and data availability statement and release timing, given
the pending patent filing; a conservative draft is ``available from the
corresponding author on reasonable request''}
